\scanpage{49}
le cas où tous les $E_i$ sont réduits à un
élément \deleted{\unsure{et}} ($\Leftrightarrow$ \deleted{\unsure{il}} un des $E_i$ l'est $\Leftrightarrow E$ l'est),
cas dit ``trivial''. Donc une $n$-alité non
triviale peut être considérée comme
un ens $E$ \unsure{munie|muni} d'une partie
$I\subset\operatorname{Quot}(E)$, avec $\operatorname{Card} I=n$, et \unsure{ou|un|satisfaisant}
l'axiome évident traduisant (1).

Autre version : Soit \deleted{\unsure{$S$}}
\[
S=\coprod_{i\in I}E_i\qquad\text{(``sommets'')}
\tag*{(2)}
\]
et \deleted{\unsure{$f$}} soit
\[
R\subset S\times E \tag*{(3)}
\]
défini par
\[
(\underset{\substack{\Vert\\(i\in I,s_i\in E_i)}}{s},
 \underset{\substack{\uparrow\\E}}{x})\in R\quad\Longleftrightarrow\quad
s=\varphi_i(x)
\]

\deleted{Alors on a ceci}

Considérons
\[
E\longrightarrow \unsure{\mathcal P|P|\mathfrak P}(S)\qquad x\longmapsto R(\unsure{a|x})\subset\deleted{E}S
\]
les éléments de \deleted{$R(x)$} la forme $R(x)$ sont
appelés $n$-uples de $S$. On dit que
\deleted{\unsure{a)} $E\longrightarrow\mathcal P(S)$ est \unsure{+}}
$a\in\mathcal P_2(S)$ est une \uline{arête} si \unsure{$a$|\deleted{$a$}|elle} est \unsure{contenu|contenue}
dans un $n$-uple, i.e. $\{s,t\}$ est une arête
sss $s\ne t$ et $\exists x\in E$ tel que $\{s,t\}\subset R(x)$ i.e.
% archive sheet number: 48
